% !TEX program = pdflatex
% !TEX encoding = UTF-8
% Standalone submission document; compile this file directly.
\documentclass[conference,a4paper]{IEEEtran}
\usepackage[T1]{fontenc}
\usepackage{xurl}
\usepackage{balance}
\usepackage[hidelinks]{hyperref}
\hypersetup{
  pdftitle={FallGuard: Yaorong Huang Individual Contribution},
  pdfauthor={Yaorong Huang},
  pdfsubject={COMP6131 Fundamentals of the Internet of Things, Group 7}
}
\urlstyle{tt}
\newenvironment{evidence}{\par\smallskip\begingroup\footnotesize\raggedright\noindent}{\par\endgroup}
\begin{document}
\title{FallGuard: Individual Work Contribution and Reflection}
\author{
\IEEEauthorblockN{Yaorong Huang}
\IEEEauthorblockA{Group 7, COMP6131 Internet of Things Essentials\\
Faculty of Applied Sciences, Macao Polytechnic University, Macao SAR, China}
}
\maketitle
\begin{abstract}
I served as group leader and the principal implementation and integration contributor to FallGuard. My work connected quality-aware temporal modeling, experiment tooling, portable model export, independent Raspberry Pi inference, persistent cloud delivery, complete operator workspaces, and the final technical report and presentation. This document records the tasks, implementation problems, solutions, and concrete deliverables from that work.
\end{abstract}
\section{Yaorong Huang}
\subsection{Role and overall responsibility}

I served as group leader and principal implementation and integration contributor. I carried out most of the project's code implementation and took primary responsibility for the technical manuscript and web presentation. My work connected the research model to a running IoT system: model design and preprocessing, experiment tooling, portable inference, Raspberry Pi deployment, local persistence, cloud synchronization, operator workspaces, and reproducibility materials. I also coordinated the interfaces between the team's data, training, cloud, and reporting tasks and completed the final integration and consistency checks.

\subsection{Model design, implementation, and experiment tooling}

I implemented MaskedBiMamba, combining temporal attention, independent forward and backward Mamba processing, training-time frame masking, and quality-weighted pooling. I defined how normalized skeleton coordinates and joint confidence enter the temporal model and how the separate quality features determine frame weights. This required implementing the input representation, branch alignment, feature fusion, pooling, and classification interface, then connecting those components to the training and evaluation code. The resulting model provides a common input/output contract for research evaluation and deployment.

I implemented and organized the model-comparison and component-ablation tooling, including configuration files, random seeds, checkpoint selection, metric calculation, and result aggregation. The evaluation covered baseline comparisons, component ablations, and a separate TCNTE reproduction track under a consistent protocol. My work covered the complete comparison pipeline, from recorded configurations and checkpoint selection to aggregated results and figures. I also organized primary robustness evaluation, external Le2i evaluation, and replay of the deployed confirmation rules.

I connected experiment outputs to the manuscript's tables and plots, checking the distinction between window predictions, video predictions, merged events, and confirmed alerts. The results showed an advantage in window-level precision and stable performance when frames were missing. These results gave the model-development work a clear evaluation outcome and a concrete rationale for its deployment.

\subsection{Model export and independent Raspberry Pi inference}

I converted the selected temporal checkpoint into a portable two-input ONNX model and maintained its input shapes, normalization, class order, model card, and checksums. I implemented the portable inference path so that the Pi classifier does not depend on the CUDA training extensions. The deployment accepts skeleton sequences and their corresponding quality features, matching the model used in the experiment pipeline.

I integrated USB-camera and video acquisition, YOLOv8s-pose extraction accelerated by Hailo, a CPU pose fallback, and local ONNX classification. I implemented timestamped sliding windows, interpolation, pose-quality checks, and handling of missing or irregular observations. I also implemented the local alert transitions and connected them to stored records. This work turned the trained classifier into a service that can acquire input, classify windows, maintain state, and preserve results on the Raspberry Pi without an external network.

I supported the paired hardware evaluation and incorporated its measurements into the deployment and report. The paired comparison demonstrated a substantial improvement in pipeline throughput with Hailo-assisted pose extraction. Temporal classification remained on the Pi CPU. The deployed result therefore reflects acceleration of the complete implemented pipeline, together with a usable CPU fallback.

\subsection{Persistent delivery and complete operator workspaces}

I integrated SQLite record storage, an independent outbox, and cloud synchronization. The design preserves record identifiers and original capture times, commits cloud records before acknowledgement, and safely retries pending records after reconnection. I connected these components to the cloud archive and the private device-status bridge. Delivery and recovery tests confirmed that records remained available locally and were synchronized after the connection recovered.

I completed the local and cloud workspaces for detection overview, history, and device diagnostics. Beyond a basic prediction display, I implemented live preview, result freshness, historical probability charts, filters, pagination, CSV export, queue status, and resource diagnostics. I added explicit handling for incomplete poses, expired results, paused cameras, and lost device connections, so that the current display follows the operating state while saved history remains available.

I verified service startup and shutdown, archive queries, authentication, delivery, and recovery using the software suite and the actual page scripts. Deployment updates preserved existing SQLite data, checked service health and source hashes, and retained rollback copies. I also restored the browser-camera demonstration alongside the offline workspace, including its landmark-to-classifier adapter and current-model inference. This provided both the deployed Pi workflow and a computer-based demonstration entry.

\subsection{Problems encountered and solutions implemented}

One integration problem was maintaining the same model input in training and deployment. I made the pose normalization and separate quality input explicit, aligned the timestamped buffer with the training window definition, and checked preprocessing and PyTorch/ONNX numerical agreement. The export checks confirmed close numerical agreement between the original and deployed models.

A second problem was keeping inference and recorded results available during network interruption. I separated camera inference, local storage, the outbox, and synchronization, then exercised restart, retry, and lost-acknowledgement behavior. UUID-based deduplication and acknowledgement after durable storage preserved the recorded history during recovery.

A third problem was that the earlier interface did not provide enough information for independent operation and troubleshooting. I replaced it with the full local/cloud workspaces and added freshness checks, history inspection, and diagnostics. I also corrected the preprocessing order in the missing-joint test path and updated the result selection, figures, and manuscript together. Resolving these issues required changes across model code, deployment services, interfaces, and documentation.

\newpage
\subsection{Manuscript, presentation, and reproducibility publication}

I took primary responsibility for writing and revising the IEEE LaTeX report: method notation and equations, architecture diagrams, experiment protocols, results, system implementation, figures, captions, and references. I revised the argument around the model's precision, frame-loss tolerance, edge throughput, and recoverable cloud delivery, while keeping the numerical statements consistent with the saved results. I also maintained the Chinese mathematical and system-workflow guides and the project brief, providing explanations for both technical review and presentation preparation.

I developed the GitHub Pages presentation, interactive comparison charts, a searchable results archive, bilingual speaker notes, and presenter and rehearsal tools. I synchronized the data, explanations, source files, and published outputs across the report, guides, presentation, deployed software, and repository.

I organized the public reproducibility release: experiment and deployment source, dependency files, generic configuration and service templates, selected weights, ONNX model, asset hashes, dataset provenance, and run instructions. I used the software test suite, export consistency checks, release manifests, and GitHub verification workflow to check the handover. I prepared the final report source, independent contribution document, presentation PDF, and fixed repository submission version as a coherent set of deliverables.

\subsection{Deliverables and reflection}

\begin{evidence}
Main deliverables: \path{code/src/fallbench/}, \path{code/configs/}, \path{code/scripts/}, \path{deployment/}, \path{artifacts/models/}, \path{webui/}, \path{tests/}, \path{results/}, \path{manuscripts/final-report/}, \path{documentation/zh/}, \path{docs/}, and the reproducibility/submission guides.
\end{evidence}

My contribution extended from algorithm implementation to the final running system and its submitted evidence. The most demanding part was keeping these layers consistent: the model input used by training, the exported classifier, the device buffer, the cloud record, and the value shown to an operator. Working through that chain strengthened my ability to debug across components and to connect a research result to a maintainable IoT implementation. I also learned to treat experiment organization, recovery behavior, technical writing, and reproducibility as substantive engineering work, with concrete outputs that can be inspected and run.

\end{document}
